\documentclass[10pt,twocolumn]{article}

% ---------------------------------------------------------------------------
% OpenThomas: The Harness Is the Edge
% A venue-neutral draft (primary target: ACM ICAIF). Compiles with pdflatex
% and only standard CTAN packages -- no conference .sty required. To retarget
% at ICAIF, swap in \documentclass{acmart}; at COLM/NeurIPS, their style file.
% ---------------------------------------------------------------------------

\usepackage[utf8]{inputenc}
\usepackage[T1]{fontenc}
\usepackage[margin=0.85in]{geometry}
\usepackage{amsmath,amssymb}
\usepackage{booktabs}
\usepackage{microtype}
\usepackage{xcolor}
\usepackage{graphicx}
\usepackage[hidelinks]{hyperref}
\usepackage{url}
\usepackage{enumitem}
\setlist{nosep,leftmargin=1.2em}

\title{\textbf{The Harness Is the Edge:\\ A Disciplined, Self-Improving Language-Model Agent\\ for Weather Prediction Markets}}
\author{Thomas Wang\\ \small \href{mailto:thomas@openthomas.com}{thomas@openthomas.com}\\ \small \url{https://github.com/PredictionMarketTrader/openthomas} \quad\textbullet\quad \url{https://openthomas.com}}
\date{\today}

\begin{document}
\maketitle

\begin{abstract}
Every frontier language model that has been given real capital to trade
prediction markets autonomously has lost money. We argue this is not a
forecasting-accuracy problem but a \emph{harness} problem: the profitable
signature in the record is selectivity, sized asymmetry, exit discipline, and
calibration --- properties of the trading system around the model, not of the
model itself. We present \textbf{OpenThomas}, an open-source autonomous agent
built to be that harness, and deliberately pointed at the one market category
where a harness can be measured cleanly: daily station-temperature markets on
Kalshi and Polymarket. Weather markets settle daily against an objective public
record (the U.S. National Weather Service climatological report for a named
station), which turns them into a \emph{strong-evaluator} domain: dense,
objective, and reproducible in leak-free replay. OpenThomas prices every strike
from a consensus of seven operational global numerical-weather-prediction (NWP)
models plus its own GraphCast point forecast and GenCast ensemble; learns each
station's systematic bias and spread from leak-free hindcasts; lets a language
model make a \emph{bounded} adjustment; blends with the market prior; and sizes
every position with fractional Kelly under hard, deterministic risk limits the
model can never override. On top of this, the system improves itself: an
evolution loop mines its own journal, proposes bounded mutations, and must clear
a kernel-owned promotion gate on leak-free replay before anything takes effect
--- with the evaluator structurally off the agent's edit surface, so
self-improvement cannot redefine ``improvement.'' In a 21-day replay against real
historical order books with fees, a naked model-vs-market strategy loses
\$0.038/contract; adding learned station bias, a market-prior blend, and
timing hygiene walks it, layer by layer, to breakeven --- before any of the
live-only advantages are switched on. We report the losing numbers as well as
the gains, and run the reference agent in public. Our contributions are (i) the
harness-is-the-edge thesis operationalized as a concrete, reproducible system;
(ii) a leak-free hindcast/replay evaluation methodology for market agents; and
(iii) a self-improvement architecture whose safety derives from keeping the
judge off the edit surface, evaluated on a domain with a genuinely strong judge.
\end{abstract}

\section{Introduction}

The obvious way to build a trading agent on top of a large language model (LLM)
is to ask the model what will happen and bet accordingly. The empirical record
says this does not work. The \emph{Prediction Arena} benchmark
\cite{predarena} gave six frontier models \$10{,}000 each and let them trade
real prediction markets autonomously for 57 days. Every one of them lost money,
between $-16\%$ and $-30.8\%$. Research \emph{volume} had no correlation with
returns; token spend had none either. What separated the least-bad runs from the
worst was not a better world model but better \emph{behavior}: being right early,
sizing up when confident and down when uncertain, holding winners to settlement
rather than exiting early, and --- above all --- not trading at all when there
was no edge. The single profitable run on record ($+10.9\%$) was selective (112
trades), asymmetric (average win \$63.89 vs.\ average loss \$3.23), and
low-drawdown (4.1\%).

That profile is a property of the \emph{harness} --- the scaffolding that
decides when the model is allowed to act, how much, and under what constraints
--- not of the model's raw forecasting ability. Our central claim follows:

\begin{quote}
\emph{The LLM is not the edge. The harness is the edge.}
\end{quote}

This reframing has a sharp methodological consequence. If the edge lives in the
harness, then to make progress we need a domain where harness quality is
\emph{measurable} --- densely, objectively, and without leakage. Most prediction
markets fail this test: they settle rarely (one data point every few months),
their resolution is ambiguous, and any backtest risks letting the replayed
market ``know'' something the replayed model did not. \textbf{Weather markets are
the exception.} A daily station-temperature market (``Will the high in Central
Park exceed 83\textdegree F today?'') resolves against the NWS climatological
report for one named station, every day, with no ambiguity. Dozens of
ground-truth data points arrive per day. This is the fastest, cleanest learning
loop in prediction markets, and it is what makes both the trading edge and,
later, recursive self-improvement, something you can put a number on.

We present \textbf{OpenThomas}, an open-source ($\approx$7{,}300 lines of Python,
182 tests) autonomous agent that operationalizes the harness thesis on weather
markets. Its design commitments are:

\begin{itemize}
\item \textbf{Selectivity over activity.} ``No trade'' is the default action.
The agent trades only when $|\,\hat{p}-\text{price}\,|$ clears an explicit
expected-value threshold after fees and spread.
\item \textbf{A statistical baseline the model may adjust but never replace.}
Every strike is priced from a seven-model NWP consensus, corrected by
per-(station, lead) bias and spread learned from leak-free hindcasts, and hard-
truncated by what today's thermometer has already locked in. The LLM is clamped
to $\pm 0.15$ around this baseline.
\item \textbf{Hard risk constraints outside the model.} Fractional Kelly sizing,
per-market / per-event / per-category caps, a fee-aware EV threshold, a
longshot-zone filter, and a drawdown kill-switch are deterministic code. The
model proposes; the risk engine disposes.
\item \textbf{Calibration as a first-class metric,} tracked per (station, lead)
as Brier skill against the market, and used to license sizing only where skill
is positive.
\item \textbf{Self-improvement as a gated loop} in which the evaluator is
structurally off the agent's edit surface (Section~\ref{sec:rsi}).
\end{itemize}

We contribute: (1) the harness-is-the-edge thesis operationalized as a concrete
system, with an honest empirical account including the losing runs; (2) a
leak-free hindcast/replay methodology for evaluating market agents where the
replayed market is never allowed to know what the replayed model does not; and
(3) a self-improvement architecture whose safety comes from a kernel/agent plane
split that keeps the judge off the edit surface, stress-tested on a domain with a
genuinely strong judge. The system, the evaluation harness, and the reference
agent's live track record are all public.

\section{Background and related work}

\paragraph{LLMs as forecasters and traders.} A growing literature evaluates LLMs
as probabilistic forecasters. The consistent finding is that \emph{pure} LLM
forecasts underperform market consensus, but a \emph{blend} of the model with the
market price can beat the market \cite{blend}. Prediction Arena \cite{predarena}
extends this from forecasting to \emph{live autonomous trading} and finds
uniform losses, isolating behavioral rather than epistemic failure modes. Our
work takes the blend result and the behavioral diagnosis as design inputs: we
blend with the market prior by construction, and we push the behavioral factors
(selectivity, asymmetry, exit discipline) into deterministic code rather than
prompting for them.

\paragraph{Neural weather prediction.} Machine-learned global weather models now
match or exceed operational physics-based systems: GraphCast \cite{graphcast}
produces deterministic 10-day forecasts, and GenCast \cite{gencast} produces
calibrated ensembles via a diffusion sampler. These open-weights models can be
run on commodity accelerators. OpenThomas uses them as \emph{additional voters}
in a forecast consensus and, crucially, post-processes them per settlement
station --- the place where public global forecasts are systematically biased and
where a small, learnable, local correction becomes a durable edge.

\paragraph{Prediction-market microstructure.} The favorite--longshot bias ---
low-priced contracts are systematically overbet --- is well documented
\cite{longshot}; we encode it as a hard longshot-zone filter. Cross-venue
listing of the same real-world event (Polymarket vs.\ Kalshi) creates episodic
arbitrage, which our scanner flags deterministically.

\paragraph{Self-improving agents.} Recent framing of near-term recursive
self-improvement (RSI) locates the leverage not in model weights but in the
\emph{harness} --- how the agent plans, acts, remembers, and is evaluated --- and
identifies \emph{evaluator strength} as the first bottleneck \cite{harness}.
Open-ended methods such as the Darwin Gödel Machine \cite{dgm} keep an archive of
past variants as mutation parents to resist diversity collapse. OpenThomas adopts
both ideas and adds a security model: because weather markets supply an unusually
strong evaluator, the risk is no longer a weak judge but a \emph{gameable} one,
so we make the judge structurally unwritable by the optimizer.

\section{Why weather markets are a strong-evaluator domain}
\label{sec:weather}

Three properties make daily station-temperature markets uniquely suited to
measuring harness quality.

\textbf{Dense, objective settlement.} Each market resolves against the NWS
Climatological Report (product type \texttt{CLI}) for one named station, once per
day. There is a single number --- the day's observed maximum or minimum --- and no
interpretive latitude. A generalist prediction-market agent might see a few
dozen settlements per \emph{year}; a weather desk on ten U.S.\ cities
(highs and lows) sees dozens per \emph{day}.

\textbf{A learnable, local edge.} Markets settle on a \emph{specific} station
(NYC $=$ Central Park; Chicago $=$ Midway), and global forecast models have
systematic, station-specific biases there. Our first 90-day hindcast found the
Miami consensus running $1.8$--$2.1$\textdegree F cold ($\sigma\,1.7$),
Philadelphia $+2.3$\textdegree F, and LAX about $1$\textdegree F warm. These are
biases you can estimate from free public data, per station, offline --- a moat
that does not depend on privileged information, only on the discipline to
measure and apply it.

\textbf{Leak-free reproducibility.} Historical ``as-of'' forecast data (what a
model predicted at a past issue time) and historical order books are both
available, so a backtest can present the replayed agent with exactly the
information it would have had --- and no more. This is what lets us evaluate not
just the forecast but the \emph{decision rule}, and later, self-improvement,
without the replay silently cheating.

The same three properties are what the RSI literature asks for and rarely gets:
an objective, dense, reproducible evaluator \cite{harness}.

\section{System architecture}
\label{sec:arch}

OpenThomas is organized as a stack of layers with a strict rule about who is in
control at each one (Table~\ref{tab:layers}). Everything that touches money ---
sizing and exposure --- is deterministic code; the language model is confined to
a bounded adjustment in the middle of the pipeline.

\begin{table*}[t]
\centering
\small
\begin{tabular}{@{}p{2.4cm}p{9.5cm}p{3.0cm}@{}}
\toprule
\textbf{Layer} & \textbf{What it does} & \textbf{Who's in control} \\
\midrule
Edge scanner & Filters markets to plausible mispricings; flags cross-venue arbitrage and longshot/resolution-trap zones & deterministic code \\
Weather baseline & Seven-model NWP consensus $+$ own GraphCast/GenCast $\to$ per-strike probability at the settlement station, with per-(station, lead) bias/$\sigma$ from leak-free hindcasts and hard truncation by today's observed extreme & deterministic code \\
Forecast engine & LLM ensemble (median of $N$ independent estimates), grounded in resolution rules and base rates, \emph{clamped to baseline $\pm 0.15$} & your choice of model \\
Calibration & Platt-scales forecasts against your own settled history; blends with the market prior & deterministic code \\
Risk engine & Fractional Kelly, per-market/event/category caps, fee-aware EV threshold, longshot filter, drawdown kill-switch & deterministic code \\
Memory & Every forecast and fill journaled to SQLite; post-settlement reflection distills bounded playbook rules & agent, human-auditable \\
\bottomrule
\end{tabular}
\caption{The OpenThomas stack. The model \emph{proposes}; the risk engine
\emph{disposes}. Only one layer is model-controlled, and it is bounded above and
below by deterministic layers.}
\label{tab:layers}
\end{table*}

\subsection{Market connectors and paper simulation}
A unified \texttt{Market}/\texttt{Order}/\texttt{Position} model spans Polymarket
(Gamma API for discovery, CLOB for books/orders) and Kalshi (REST $+$ WebSocket,
RSA-signed, CFTC-regulated). A \emph{paper} connector wraps either venue's real
market data but simulates fills locally at bid/ask (never mid), with
mark-to-market at bid prices (liquidation value, not hope). Paper mode is the
default and needs no exchange keys, which makes the whole system runnable ---
and reproducible --- by anyone.

\subsection{The weather baseline}
For a market on station $s$, valid day $d$, kind $k\in\{\text{high},\text{low}\}$,
the baseline forms a consensus temperature $\bar{T}$ across the seven operational
global models (GFS, ECMWF-IFS, ICON, GEM, UKMO, M\'et\'eo-France, JMA) plus the
system's own GraphCast point forecast, with a cross-model spread $\sigma_{\text{model}}$.
It then applies two corrections learned strictly out-of-window:
\[
\hat{T} = \bar{T} + b_{s,k,\ell}, \qquad
\sigma^2 = \sigma_{\text{model}}^2 + \sigma_{s,k,\ell}^2 ,
\]
where $\ell$ is the forecast lead, $b_{s,k,\ell}$ is the learned station bias and
$\sigma_{s,k,\ell}$ its residual spread. GenCast contributes a
\emph{flow-dependent} multiplier on $\sigma$ (an uncertain synoptic setup widens
the distribution) rather than a point estimate, because its 12-hour steps
undersample the diurnal peak. The per-strike probability is then the Gaussian (or
empirical) mass on the correct side of the strike, with one more hard step:
\emph{truncation}. If today's thermometer has already recorded a value that makes
one side of the strike impossible (the observed high already exceeds the floor),
the probability is clamped accordingly. This is where a naive backtest leaks, and
where Section~\ref{sec:eval} spends most of its care.

\subsection{Forecast engine: bounded LLM adjustment}
The language model's role is deliberately small. It receives the baseline
probability, the resolution rules, base rates, and a news brief, and returns an
adjusted probability --- which is then \emph{clamped to baseline $\pm 0.15$}. The
statistics do the heavy lifting; the model is there for what the statistics
cannot see (a blown front, a stale consensus, a discontinuity in the question
text), not to substitute its own guess for a calibrated ensemble. Because of the
clamp, mid-size local reasoning models suffice, which keeps the agent runnable on
self-hosted open-weights models with no API bill and full privacy. The engine
takes the median of $N$ independent estimates under different framings to reduce
single-sample variance.

\subsection{Calibration and market-prior blend}
Once $\geq 30$ markets settle in a category, OpenThomas fits Platt scaling to
correct the model's systematic over/under-confidence, then blends the calibrated
forecast with the market price. The blend is not a hedge --- it is the empirically
supported position that a model-market blend beats either alone \cite{blend}. The
blend weight is an entry-\emph{selectivity} parameter (how much to trust our view
vs.\ the crowd), and as such is a candidate for self-improvement under kernel
bounds (Section~\ref{sec:rsi}); the sizing that follows is not.

\subsection{Risk engine (deterministic)}
The risk engine takes bankroll, risk profile, the (blended) probability and
confidence, and market price/liquidity, and returns an approved size or a
rejection naming the binding constraint. It applies fractional Kelly (profile-
scaled $0.15\times$--$0.33\times$), per-market concentration caps, per-category
and per-event correlation caps (correlated settlements produced Prediction
Arena's largest single-session losses), a fee-aware EV threshold, a longshot-zone
filter, a solvency check including fees, and a drawdown kill-switch that halts
trading and requires a human to resume. None of these are reachable by the model
or by the self-improvement loop; a kernel policy test fails any change that tries
to move a safety rail into the optimizer's reach.

\section{Self-improvement with the judge off the edit surface}
\label{sec:rsi}

OpenThomas updates OpenThomas. The operator (a human, or a coding assistant) is
moved \emph{up} the stack --- setting bounds and reviewing lineage --- and is
forbidden from hand-tuning trading parameters; that is what the evolution loop is
for. The design problem is that self-improvement is safe only if improvement
cannot redefine ``improvement.'' An optimizer that can edit its own judge will
find it easier to lower the bar than to clear it.

\paragraph{Kernel/agent planes.} The codebase is split into two planes, and the
split \emph{is} the security model (Table~\ref{tab:planes}). The \textbf{kernel}
plane --- promotion gate, parameter bounds, the risk engine and all safety rails,
the truth pipeline, the journal schema --- is written only by the operator; the
evolution loop has no write path to it. The \textbf{agent} plane --- decision-rule
parameters, the forecast prompt template, and (on the roadmap) workflow
structure, strategy code, and the evolver itself --- is written by the evolution
loop, but only \emph{through} the kernel gate. The line is drawn precisely:
entry-selectivity knobs (minimum edge, market-prior blend weight) are genome
parameters under kernel bounds, while sizing, exposure caps, and the kill-switch
can never enter the parameter space.

\begin{table}[t]
\centering
\small
\begin{tabular}{@{}p{1.4cm}p{4.2cm}p{1.7cm}@{}}
\toprule
\textbf{Plane} & \textbf{Contents} & \textbf{Written by} \\
\midrule
kernel & bounds, promotion gate, risk engine \& safety rails, truth pipeline, journal schema & operator only \\
agent & decision-rule params, forecast prompt template; later: workflow, strategy code, the evolver & evolution loop, \emph{via the gate} \\
\bottomrule
\end{tabular}
\caption{The plane split is the security model: the judge is not on the
optimizer's edit surface.}
\label{tab:planes}
\end{table}

\paragraph{Three nested loops.} A fast \emph{trading} loop (minutes) runs the
champion parameters and never waits on anything below it. A \emph{reflection}
loop (on settlements) distills bounded playbook rules. An \emph{evolution} loop
mines failure evidence from the journal, proposes mutations --- LLM-directed by
the evidence, plus Gaussian mutations drawn from an \emph{archive} of past
variants (Darwin Gödel Machine-style, against diversity collapse
\cite{dgm}) --- and submits them to the gate.

\paragraph{The gate (kernel, frozen).} A candidate is a \emph{strategy callable}
scored on leak-free replay with four anti-gaming properties:
\emph{held-in/held-out} splits (winners are picked on older days; the freshest
settled days only get a veto, and the window rolls forward daily so repeated
meta-cycles face new data); \emph{totals-not-ratios} scoring with a minimum trade
count (a rule that trades twice cannot win on ratio; a rule that trades nothing
cannot win on drawdown); \emph{out-of-sample regret rollback} (each cycle first
re-scores the active generation against its parent on the fresh window and
reverts if the parent now wins); and a \emph{Brier veto} for the forecast
operator (a prompt that wins PnL while worsening calibration got lucky, and luck
does not promote). Generation 0 is the operator's config and applies no
overrides --- the human's settings win until the loop has promoted
evidence-backed changes past them.

\paragraph{One gate, growing write surface.} Improvement axes are unified as
mutation operators feeding the \emph{same} accept mechanism: playbook rules and
calibration refits (shipped), decision-parameter mutation and forecast-prompt
mutation (shipped), then workflow structure, then agent-plane code diffs (which
add shadow paper-trading and human diff review), then LoRA/GRPO weight updates
with a time-discounted reward $R\cdot e^{-\lambda\,\text{days-to-close}}$ --- a
trained adapter is just another candidate at the same gate. The admission rule is
strict: a parameter enters the space only when the gate can actually discriminate
on it. \emph{Tuning what you cannot measure is drift, not improvement.} This is
why code diffs come last --- they need the strongest evaluation (replay \emph{plus}
shadow trading \emph{plus} human review of the diff) --- and why the whole
enterprise is anchored to a domain with a strong evaluator in the first place.

\section{Leak-free evaluation methodology}
\label{sec:eval}

The hardest part of evaluating a market agent is not computing PnL; it is not
lying to yourself. Two commands encode the discipline.

\texttt{hindcast} teaches the baseline each station's bias from historical as-of
forecasts and official settlements, strictly \emph{before} any trading window it
will be evaluated on. \texttt{replay} then backtests the \emph{decision rule}
against real historical order books, drawing prices from snapshot candlesticks
that \emph{precede} the temperature extreme being bet on, with station bias/$\sigma$
learned only out-of-window. The governing rule, learned the hard way when an
early backtest snapshotted the dawn low after it had already been realized:
\emph{never let the replayed market know something the replayed model does not.}
A single timing leak turned a losing strategy into a fake winner; fixing it
turned it back.

This methodology is what makes the self-improvement gate trustworthy: the same
\texttt{collect\_rows} pipeline that scores a candidate generation is the one that
enforces no-peeking, and it lives in the kernel where the optimizer cannot touch
it.

\section{Empirical results}
\label{sec:results}

We report an ablation from a 21-day replay on real Kalshi order books with fees
included (Table~\ref{tab:ablation}). We stress that these are small-sample,
paper-mode numbers on a baseline \emph{deliberately stripped} of every live-only
advantage (same-day model runs, intraday observations, LLM adjustment, and the
calibration layer); they are an existence proof about \emph{where the edge comes
from}, not a claim of profitability. The live agent trades with all four
advantages on, and whether they clear the bar is precisely what the public paper
run exists to answer.

\begin{table}[t]
\centering
\small
\begin{tabular}{@{}p{5.2cm}r@{}}
\toprule
\textbf{Strategy (21-day replay, fees in)} & \textbf{PnL/contract} \\
\midrule
Naked model-vs-market & $-\$0.038$ \\
\quad $+$ learned station bias $+$ market-prior blend & $-\$0.018$ \\
\quad\quad $+$ timing-leak fix (low snapshot at 06Z) & $\approx \$0.000$ \\
\bottomrule
\end{tabular}
\caption{Each layer of discipline is worth $\approx 2$\textcent/contract. The
naked strategy loses because the market has already priced the morning
observation --- the harness thesis, confirmed in miniature. Breakeven is reached
\emph{before} any live-only advantage is counted.}
\label{tab:ablation}
\end{table}

Two findings are worth drawing out. First, the naked model-vs-market strategy
\emph{loses} ($-\$0.038$/contract over 417 trades), because by late morning the
order book has already digested the day's observations --- exactly the ``the
model is not the edge'' prediction, borne out on real books. Second, the
station-bias signal is real and sizable: a 90-day hindcast recovered stable,
station-specific biases (Miami $+1.8$--$2.1$\textdegree F, Philadelphia $+2.3$,
Atlanta $+2.0$, LAX $-1$) that persist out-of-sample. Discipline --- learned bias,
a market blend, and timing hygiene --- is what converts a losing naive strategy
into a breakeven one; the ablation attributes roughly two cents per contract to
each layer.

We deliberately publish the negative and near-zero numbers. The honest claim is
narrow and, we think, more useful than a headline return: on weather markets the
structural edges are \emph{real but small}, discipline is what preserves them, and
this is the market category where that statement can be checked.

\section{Discussion and limitations}

\paragraph{Small samples, paper mode.} The replay is 21 days and simulated at
real bid/ask; it is an ablation, not a track record. The reference agent's live
paper run is the ongoing, public test, reported with wins and losses alike.

\paragraph{Capacity is the real ceiling.} Daily-temperature markets are thin.
A survey of seven Kalshi cities found on the order of a few hundred thousand
dollars of daily notional in aggregate, which bounds extraction at low-single-
digit millions per year even under optimistic assumptions. This is a research and
methodology contribution about \emph{measurable} edge, not a claim of scalable
returns.

\paragraph{Replay fidelity.} Replay has no news brief, no intraday observations,
no forecaster-discussion text, and no calibration layer; it measures the
template's skill at adjusting the statistical baseline, which is the dominant live
pathway for weather markets, but it is not the whole live agent. A residual leak
risk --- an LLM ``remembering'' recent weather from pre-training --- is negligible
with a local model whose cutoff precedes the window, but must be rechecked
whenever the forecaster is swapped.

\paragraph{Generalization.} The strong-evaluator property is what makes weather
the right first domain; it is also what most other markets lack. Extending the
harness (and the self-improvement gate) to weaker-evaluator markets is future
work, and the journal already archives each live forecast's inputs so that a
future trace-replay evaluator can begin to close that gap.

\section{Reproducibility, ethics, and broader impact}

The system, the evaluation harness, and the reference agent's live feed are
public; paper mode requires no exchange keys and reproduces the replay numbers
from public data. Live trading is off by default and requires two independent
explicit opt-ins; the agent has no deposit, withdrawal, or transfer capability
and can only trade within an allocated bankroll. We ship hard position limits and
a drawdown kill-switch, and we state plainly that no software makes trading safe:
roughly 70\% of Polymarket addresses have lost money, and every layer of this
paper is about losing less, not winning easily. On the self-improvement side, the
safety argument is architectural rather than behavioral: the judge is off the
optimizer's edit surface by construction, safety rails are outside the parameter
space by a kernel policy test, and every future promotion is an auditable commit
authored by the loop itself.

\section{Conclusion}

Autonomous LLM traders lose money not because they forecast badly but because
they lack discipline. OpenThomas puts the discipline in deterministic code, uses
a language model only for a bounded adjustment it is well-suited to, and points
the whole system at weather markets --- the one category where the resulting edge,
and the system's ongoing self-improvement, can be measured densely, objectively,
and without leakage. The edge we find is real, small, and preserved by discipline.
We report it in public, losses included, and we invite the community to check it.

\small
\begin{thebibliography}{9}
\bibitem{predarena} Prediction Arena: evaluating frontier language models as
autonomous prediction-market traders. arXiv:2604.07355, 2026.
\bibitem{blend} Calibrated LLM forecasts and market-prior blending.
arXiv:2511.07678, 2025.
\bibitem{graphcast} R.~Lam et al. GraphCast: Learning skillful medium-range
global weather forecasting. \emph{Science}, 2023.
\bibitem{gencast} I.~Price et al. GenCast: Probabilistic weather forecasting with
machine learning (diffusion ensembles). \emph{Nature}, 2025.
\bibitem{longshot} J.~Whelan et al. The favorite--longshot bias in prediction
markets. SSRN 5502658, 2026.
\bibitem{harness} L.~Weng. Harness Engineering for Self-Improvement.
\url{https://lilianweng.github.io/posts/2026-07-04-harness/}, 2026.
\bibitem{dgm} The Darwin G\"odel Machine: open-ended self-improvement with a
variant archive. 2025.
\end{thebibliography}

\end{document}
